% Volume VIII: Applications and Case Structures
% Part XVI. Failure Modes of the Theory
% Chapter 95: The Abuse of “More Context”
% Spine: proof-spines/ch095-spine.md

\chapter{The Abuse of ``More Context''}
\label{ch:ch095}

Requests for context can be legitimate, but they can also indefinitely postpone recognition of clear evidence. This chapter defines diminishing contextual returns and asks when further expansion no longer materially changes the conclusion relevant to the decision at hand. Its subject is the recurrent maneuver by which every supplied clarification produces a new demand for background, sequence, motive, culture, or comparative framing.

The point is not to defend decontextualized judgment. It is to distinguish context that discriminates among live explanations from context that merely delays accountability until witnesses scatter and attention fades. The chapter therefore offers a stopping rule: inquiry should keep widening only while the added frame could still alter the warranted threshold judgment, not simply because more surrounding information can always be imagined.
